% Compilar con:  latexmk -xelatex 03-arquitectura.tex
\documentclass[11pt,a4paper]{article}

\usepackage{fontspec}
\usepackage[spanish,es-noshorthands]{babel}
\usepackage{geometry}
\geometry{margin=2.4cm}
\usepackage{xcolor}
\usepackage{array}
\usepackage{booktabs}
\usepackage{tabularx}
\usepackage{enumitem}
\usepackage{titlesec}
\usepackage{fancyhdr}
\usepackage[most]{tcolorbox}
\usepackage{newunicodechar}
\usepackage{hyperref}
\usepackage{fvextra}

% Fuente monoespaciada con soporte de caracteres box-drawing
\setmonofont{DejaVu Sans Mono}[Scale=0.82]

% Entorno para bloques de código (verbatim con font Unicode)
\DefineVerbatimEnvironment{Diagram}{Verbatim}{%
  fontfamily=tt, fontsize=\small,
  breaklines=false, commandchars=\\\{\}}

% --- Símbolos Unicode ---
\newunicodechar{→}{\ensuremath{\rightarrow}}
\newunicodechar{←}{\ensuremath{\leftarrow}}
\newunicodechar{≤}{\ensuremath{\leq}}
\newunicodechar{≥}{\ensuremath{\geq}}
\newunicodechar{×}{\ensuremath{\times}}
\newunicodechar{≈}{\ensuremath{\approx}}
\newunicodechar{·}{\textperiodcentered}
\newunicodechar{†}{\textsuperscript{\dag}}

% --- Colores y estilo (coherentes con el resto del paquete) ---
\definecolor{accent}{HTML}{0F6E8C}
\definecolor{accentdark}{HTML}{0A4B60}
\definecolor{soft}{HTML}{F2F6F8}
\definecolor{diagbg}{HTML}{F8FAFB}

\titleformat{\section}{\Large\bfseries\color{accentdark}}{\thesection}{0.6em}{}
\titleformat{\subsection}{\large\bfseries\color{accent}}{\thesubsection}{0.6em}{}
\titleformat{\subsubsection}{\normalsize\bfseries\color{accent}}{\thesubsubsection}{0.5em}{}
\titlespacing*{\section}{0pt}{1.4em}{0.5em}
\titlespacing*{\subsection}{0pt}{1.0em}{0.4em}

\setlist[itemize]{leftmargin=1.3em,itemsep=2pt,topsep=3pt}
\setlist[enumerate]{leftmargin=1.6em,itemsep=2pt,topsep=3pt}

\renewcommand{\arraystretch}{1.25}
\setlength{\emergencystretch}{3em}
\setlength{\tabcolsep}{5pt}

\hypersetup{
  colorlinks=true, linkcolor=accent, urlcolor=accent, citecolor=accent,
  pdftitle={03 · Arquitectura — Starlink-Primary Resilient Gateway Lab},
  pdfauthor={Capstone de internship · Digi International}
}

% Cajas
\newtcolorbox{nota}{colback=soft,colframe=accent!35,boxrule=0.5pt,arc=2pt,
  left=8pt,right=8pt,top=5pt,bottom=5pt}
\newtcolorbox{clave}{colback=accent!6,colframe=accent,boxrule=1pt,arc=2pt,
  left=10pt,right=10pt,top=7pt,bottom=7pt}
\newtcolorbox{diagbox}{colback=diagbg,colframe=accent!25,boxrule=0.4pt,arc=2pt,
  left=6pt,right=6pt,top=6pt,bottom=6pt}

% Encabezado / pie
\pagestyle{fancy}
\fancyhf{}
\fancyhead[L]{\small\color{accentdark}03 · Arquitectura}
\fancyhead[R]{\small\color{accentdark}Starlink-Primary Gateway Lab}
\fancyfoot[C]{\small\thepage}
\renewcommand{\headrulewidth}{0.4pt}
\renewcommand{\footrulewidth}{0pt}

\newcommand{\hcell}[1]{\textbf{#1}}
\newcommand{\file}[1]{\texttt{#1}}

% ─────────────────────────────────────────────────────────────────
\begin{document}

\begin{center}
  {\Huge\bfseries\color{accentdark} Arquitectura}\\[6pt]
  {\large Estructura de cómo quedaría el laboratorio}\\[2pt]
  {\normalsize Digi Starlink-Primary Resilient Gateway Fleet}\\[10pt]
  {\small Capstone de internship · Digi International \quad|\quad Julio 2026}
\end{center}

\vspace{4pt}
\begin{nota}
\small
Este documento describe \textbf{cómo se ve el laboratorio montado}: topología, cableado y
separación por sitio, direccionamiento y plan de AS, protocolos, y flujo de datos.
Es la referencia visual y de contratos; los pasos para llegar aquí están en
\file{04-pasos-montaje.md}.

\medskip
\textbf{Convención:} un sitio en el sprint (Anillo 0), tres sitios idénticos en la expansión
(Anillo 1/2). Todo el hardware base es idéntico entre sitios; solo cambia el
\textbf{perfil} (\textit{native / bonding / candidate}) que se rota entre unidades.
\end{nota}

\vspace{2pt}
{\small\tableofcontents}
\vspace{6pt}

% ─────────────────────────────────────────────────────────────────
\section{Vista de conjunto (dos planos + control central)}

\begin{diagbox}
\begin{Verbatim}[fontsize=\footnotesize]
                     ┌───────────────────────────────────────────────┐
                     │         CONTROL CENTRAL (nube + host)          │
                     │                                                │
Internet ───────────►│  • Digi Remote Manager Premier (cloud)         │
                     │  • Core / route-collector  AS65000             │
                     │  • Terminación de overlay (IKEv2/IPsec)        │
                     │  • SIEM / receptor de Push Monitor (webhook)   │
                     │  • Plataforma de impairment (loss/jitter/lat.) │
                     │  • Servidor externo Digi WAN Bonding (Tier 3)  │
                     │  • [Fase 2] Lighthouse Enhance VM              │
                     └─────────────▲──────────────────▲──────────────┘
                                   │ overlay + BGP     │ gestión (HTTPS/API)
     ┌─────────────────────────────┼───────────────────┼─────────────────┐
     │                             │                   │                 │
┌────┴─────────┐           ┌───────┴──────┐     ┌──────┴───────┐
│  SITIO A     │           │  SITIO B     │     │  SITIO C     │
│  AS65001     │           │  AS65002     │     │  AS65003     │
│  perfil:     │           │  perfil:     │     │  perfil:     │
│  native      │           │  bonding     │     │  candidate   │
└──────────────┘           └──────────────┘     └──────────────┘
  (Anillo 0: solo Sitio A en el sprint; B y C en expansión)
\end{Verbatim}
\end{diagbox}

\medskip
Dos planos deliberadamente separados:

\begin{itemize}
  \item \textbf{Plano in-band (producción):} Starlink + 5G del IX40 → overlay → BGP → servicio.
        Es el que falla en las pruebas.
  \item \textbf{Plano out-of-band (rescate, fase 2):} Opengear OM1304 con
        \textbf{5G de otro carrier y otra energía} → Lighthouse. Debe seguir accesible cuando el
        in-band cae. Si compartieran energía o carrier, el supuesto de independencia sería falso.
\end{itemize}

% ─────────────────────────────────────────────────────────────────
\section{Anatomía de un sitio}

\begin{diagbox}
\begin{Verbatim}[fontsize=\footnotesize]
                   ┌───────────── SITIO (AS6500x) ─────────────────────┐
                   │                                                    │
Starlink           │  Starlink dishy ──(WAN primaria)──► IX40 WAN/ETH1 │
Performance Kit    │                                       (cobre)      │
(Gen 3)            │                                                    │
                   │  IX40 módem 5G ──(WAN backup in-band)─────────────►│
Carrier A (5G) ───►│                                                    │
                   │  IX40-05 ──SFP/ETH5 (1 Gb/s, SFP NO SFP+)──►      │
                   │     │        fabric de servicio (red anunciada     │
                   │     │        por BGP; IPv4 /24 + IPv6 /64)         │
                   │     │                                              │
                   │     ├── mgmt Ethernet ──► switch de gestión        │
                   │     │                                              │
                   │     ├── Digi Container (collector de métricas,     │
                   │     │     límites de recursos)                     │
                   │     │                                              │
                   │     └── DB9 serial ──► [Fase 2] OM1304 (Login)    │
                   │                                                    │
                   │  XBee Hive ──► LAN sitio ; XBee 3 Zigbee ──► sensor│
                   │     (temperatura / puerta / entrada digital)       │
                   │                                                    │
                   │  ConnectCore MP255 (sensor embebido)               │
                   │     • imagen DEY (Yocto); secure boot: fase 2      │
                   │     • agente → SIEM (buffer offline + replay)      │
                   │     • OTA vía ConnectCore Cloud Services (DRM)     │
                   │                                                    │
                   │  [Fase 2] OM1304-4E-C-LSP (OOB):                  │
                   │     • 5G Carrier B (independiente)                 │
                   │     • switch de gestión 4 puertos                  │
                   │     • consola serial al IX40                       │
                   │                                                    │
                   │  Energía: PDU/UPS primaria (Starlink, IX40, carga) │
                   │           PDU/UPS secundaria (OOB + su módem)      │
                   └────────────────────────────────────────────────────┘
\end{Verbatim}
\end{diagbox}

% ─────────────────────────────────────────────────────────────────
\section{Cableado y separación (paso a paso, por sitio)}

\begin{enumerate}
  \item \textbf{Starlink Performance Kit → IX40 \file{WAN/ETH1} (cobre).}
        \file{SFP/ETH1} queda \textbf{vacío}: al poblarlo se desactiva el puerto de cobre
        (son mutuamente excluyentes). Ver G6 para el SFP validado del segundo socket.

  \item \textbf{IX40 módem 5G → Carrier A.} WAN de respaldo in-band, con SIM/APN propio.

  \item \textbf{IX40 \file{SFP/ETH5} → fabric/carga de gateway (fibra 1\,Gb/s).}
        Aquí vive la red de servicio que el sitio anuncia por BGP.
        \textbf{Solo SFP, nunca SFP+} (límite deliberado del IX40).
        \file{SFP/ETH5} \textbf{sí} es independiente del par cobre/SFP de ETH1.

  \item \textbf{[Fase 2] OM1304 5G → Carrier B}, distinto del IX40, con antenas y
        alimentación separadas.

  \item \textbf{[Fase 2] OM1304 serial → DB9 del IX40.} El puerto IX40 en modo \file{Login};
        el IX40 es DTE RS-232 y el adaptador Opengear para X2/DTE es \textbf{\texttt{319015}}
        (DB9F–RJ45 crossover). Validar baud rate, 8N1 y flow control en banco.

  \item \textbf{[Fase 2] OM1304 Ethernet management switch → management del IX40, Hive y PDU.}
        Este enlace \textbf{no} transporta producción.

  \item \textbf{Hive → LAN del sitio; XBee 3 Zigbee → sensor físico.}
        El sensor mide temperatura, estado de puerta/alimentación o una entrada digital;
        \textbf{no} forma parte del forwarding crítico.

  \item \textbf{ConnectCore MP255 → LAN de gestión del sitio} (Ethernet).
        Recolecta red local + ambientales y envía al SIEM; recibe OTA por CCCS.

  \item \textbf{Energía:} Starlink, IX40 y carga en PDU/UPS \textbf{primaria};
        Opengear y su módem en circuito/UPS \textbf{secundario}.
\end{enumerate}

% ─────────────────────────────────────────────────────────────────
\section{Direccionamiento y plan de routing}

\subsection{Underlay (medido, nunca asumido)}

\begin{tabularx}{\linewidth}{@{}>{\bfseries}p{2.8cm}X>{\raggedright\arraybackslash}p{4.8cm}@{}}
\toprule
\hcell{Underlay} & \hcell{Detalle} & \hcell{A medir / registrar} \\
\midrule
Starlink &
  DHCPv4; donde se entregue, SLAAC (\texttt{/64}) + DHCPv6-PD (\texttt{/56}).
  CGNAT \texttt{100.64.0.0/10} por defecto; IPv4 pública opcional en planes Priority. &
  CGNAT vs.\ pública, IPv6-PD observado, MTU, latencia, jitter (G9). \\
5G IX40 (Carrier A) &
  SIM/APN propio. &
  IPv4/IPv6, MTU, NAT, soporte 5G SA/APN (G10). \\
{[Fase 2]} 5G OM1304 (Carrier B) &
  Otro operador, otro plano de energía. &
  Independencia real respecto a Starlink y Carrier A. \\
\bottomrule
\end{tabularx}

\subsection{Overlay y BGP}

\begin{diagbox}
\begin{Verbatim}[fontsize=\footnotesize]
      AS65000  (core / route-collector, endpoint público del overlay)
         ^           ^           ^
 eBGP   |           |           |   eBGP (sobre overlay IKEv2/IPsec)
         |           |           |
    AS65001      AS65002      AS65003
    Sitio A      Sitio B      Sitio C
    /24 v4       /24 v4       /24 v4     <- cada sitio origina 1 red IPv4 RFC1918
    /64 v6       /64 v6       /64 v6     <- y 1 IPv6 de documentación (2001:db8::/32)
\end{Verbatim}
\end{diagbox}

\begin{itemize}
  \item \textbf{Overlay} iniciado por el sitio hacia un endpoint central público.
        Baseline: \textbf{IKEv2/IPsec route-based con NAT traversal}
        (obligatorio por el CGNAT de Starlink). DMVPN es variante documentada por DAL,
        pero no baseline hasta probarla tras CGNAT.
  \item \textbf{eBGP} sobre el overlay entre cada IX40 y el core/collector.
  \item \textbf{Políticas:} prefix-lists de entrada/salida, límite máximo de prefijos,
        prohibición de anunciar default salvo prueba explícita, sin redistribución automática
        entre BGP y rutas de administración.
  \item \textbf{IPv6 de documentación} \texttt{2001:db8::/32}: jamás se anuncia a Internet.
  \item El core AS65000 es \textbf{estándar y vendor-neutral} porque Digi no vende un router
        de backbone carrier.
  \item \textbf{Pregunta abierta que el lab mide, no declara:} si el cambio Starlink→5G
        mantiene el túnel/BGP o provoca reconvergencia.
\end{itemize}

\subsection{Failover nativo (SureLink) — árbol DAL}

\begin{diagbox}
\begin{Verbatim}[fontsize=\footnotesize]
network interface {ETH1|Modem}
  |- ipv4 metric      (ETH1=1 menor -> mayor prioridad ; Modem=3)
  |- ipv4/ipv6 weight (métricas iguales + pesos -> ECMP; NO es bonding)
  `- surelink
       |- enable true
       |- tests   -> interface_up · ping(gateway) · ping(externo)
       |            · dns · http/https · variantes IPv6 · custom
       `- actions -> subir métrica · reiniciar interfaz · reset módem
                    · cambiar SIM · reboot · custom · power-cycle módem
\end{Verbatim}
\end{diagbox}

Reglas de diseño:
\textbf{múltiples destinos} (evitar que el monitor sea un nuevo punto único de fallo),
\textbf{acciones escalonadas} (primero métrica/reinicio; reset de módem/cambio de SIM solo tras
fallos persistentes), \textbf{histéresis + hold-down} contra flapping, y
\textbf{failback automático a Starlink solo tras una ventana estable}.

\subsection{WAN Bonding (experimento A/B)}

\begin{itemize}
  \item UI \texttt{Network > SD-WAN > WAN bonding}; CLI \texttt{network sdwan wan\_bonding};
        estado \texttt{show wan-bonding}.
  \item Requiere licencia habilitada en DRM + \textbf{servidor externo/VPS de bonding}
        + credenciales de túnel.
  \item Interfaces \texttt{ETH1} y \texttt{Modem} dentro del túnel; cifrado habilitado.
  \item Tier 3 de \textbf{1\,Gb/s} para que la licencia no sea el cuello de botella.
  \item Perfiles a comparar si la versión los expone: Automatic, Low Latency,
        Cellular Optimized.
  \item \textbf{Es una prueba separada} del failover por métrica/SureLink,
        con el \textbf{mismo} impairment trace.
\end{itemize}

% ─────────────────────────────────────────────────────────────────
\section{Plano de gestión unificado}

\begin{clave}
\small
El punto clave: \textbf{un solo plano de gestión} (DRM) opera routers, containers,
telemetría RF \textbf{y} el OTA del sensor embebido. No hay tres consolas distintas.
\end{clave}

\medskip
\begin{diagbox}
\begin{Verbatim}[fontsize=\footnotesize]
        ┌──────────────── Digi Remote Manager Premier ───────────────┐
        │                                                            │
IX40    │  • Configuration Manager (golden config + settings/sitio)  │
 ──────►│  • Alerts (Device Offline, noncompliance)                  │
        │  • Push Monitor (webhook http, topic alert_status) -> SIEM │
XBee    │  • Data streams / DataPoints (telemetría RF)               │
Hive ──►│  • Digi Containers (despliegue del collector al IX40)      │
        │                                                            │
Conn.   │  • ConnectCore Cloud Services (CCCS):                      │
Core ──►│      OTA .swu <- firmware repository de Remote Manager     │
MP255   │      (mismo tenant que gestiona los routers)               │
        └────────────────────────────────────────────────────────────┘
\end{Verbatim}
\end{diagbox}

\medskip
\subsubsection*{API — contratos base (validados)}

\begin{tabularx}{\linewidth}{@{}>{\bfseries}p{2.4cm}Xp{4.6cm}@{}}
\toprule
\hcell{Sistema} & \hcell{Base} & \hcell{Auth} \\
\midrule
DRM &
  \texttt{https://remotemanager.digi.com} · REST v1 \texttt{/ws/v1/*}
  (heredado XML/SCI \texttt{/ws/*} solo si no hay v1) &
  Basic, o API key con cabeceras \texttt{X-API-KEY-ID} /
  \texttt{X-API-KEY-SECRET}
  (\texttt{POST /ws/v1/api\_keys/inventory}). \\
Lighthouse (fase 2) &
  \texttt{https://\{host\}/api/v3.7} &
  \texttt{POST /sessions} → \texttt{session};
  luego \texttt{Authorization: Token \{session\}}
  (no \texttt{Bearer}). \\
\bottomrule
\end{tabularx}

% ─────────────────────────────────────────────────────────────────
\section{Consolidación de funciones (qué corre dónde y por qué)}

\begin{tabularx}{\linewidth}{@{}>{\bfseries}p{3.0cm}p{3.0cm}Xp{3.2cm}@{}}
\toprule
\hcell{Función} & \hcell{Dónde corre} & \hcell{Ventaja} & \hcell{Riesgo / mitigación} \\
\midrule
Collector de métricas de red &
  Digi Container en el IX40 &
  Menos hardware; despliegue central desde DRM; demuestra Containers. &
  Comparte CPU/RAM con el router → \textbf{límites de recursos + prueba de carga} para probar
  que no degrada el forwarding. \\
\addlinespace
Adquisición RF / lógica de sensores &
  XBee Hive &
  Aísla la app del forwarding; añade radio XBee. &
  Otro equipo/licencia/superficie; aceptable por el aislamiento. \\
\addlinespace
Sensor embebido / detección de anomalías &
  ConnectCore MP255 &
  Carril de software embebido (Yocto en Ring 0; secure boot/OTA en fase 2); diferenciador
  Starlink. &
  Es el track más difícil; ver \file{05-embebido-connectcore.md} y riesgos R-E*. \\
\addlinespace
OOB / rescate &
  Opengear OM1304 (fase 2) &
  Independiente en carrier y energía. &
  Se mantiene \textbf{austero} a propósito: el equipo que rescata la red no debe cargar
  experimentos. \\
\bottomrule
\end{tabularx}

% ─────────────────────────────────────────────────────────────────
\section{Topología de tres sitios y rotación de perfiles (expansión)}

\begin{tabularx}{\linewidth}{@{}>{\bfseries}p{3.2cm}Xp{4.2cm}@{}}
\toprule
\hcell{Zona} & \hcell{Componentes} & \hcell{Función / perfil} \\
\midrule
Control central &
  DRM Premier, core/collector AS65000, SIEM, impairment, servidor WAN Bonding,
  [fase 2] Lighthouse Enhance &
  Gestión de ambos planos, terminación de overlay, telemetría, alarmas, pruebas. \\
\addlinespace
Sitio A — AS65001 &
  Starlink → IX40+5G; [fase 2] OM1304; Hive+XBee; MP255; carga tras SFP &
  Perfil \textbf{native} (control): failover por métricas + SureLink. \\
\addlinespace
Sitio B — AS65002 &
  Igual que A &
  Perfil \textbf{bonding}: Starlink y 5G dentro del túnel agregado.
  \textbf{OOB con Digi Connect IT 4} (en lugar del OM1304). \\
\addlinespace
Sitio C — AS65003 &
  Igual que A; MP255 como carril OEM opcional &
  Perfil \textbf{candidate/chaos}: Containers, config nueva, firmware candidato. \\
\bottomrule
\end{tabularx}

\medskip
\begin{nota}
\small
Tras la primera campaña, los \textbf{perfiles A/B/C se rotan entre las tres unidades físicas}
para eliminar el sesgo de una unidad, un terminal Starlink o una celda concreta.
\end{nota}

% ─────────────────────────────────────────────────────────────────
\section{Resumen de contratos de direccionamiento}

\begin{tabularx}{\linewidth}{@{}>{\bfseries}p{3.4cm}p{4.8cm}X@{}}
\toprule
\hcell{Elemento} & \hcell{Valor} & \hcell{Nota} \\
\midrule
AS central      & \textbf{65000}                        & Core/collector vendor-neutral. \\
AS sitios       & \textbf{65001 / 65002 / 65003}        & Uno por sitio. \\
Red IPv4 / sitio & RFC1918 \texttt{/24}                 & Detrás de \texttt{SFP/ETH5}. \\
Red IPv6 / sitio & Documentación \texttt{/64} dentro de \texttt{2001:db8::/32} &
                   Nunca anunciada a Internet. \\
Overlay         & IKEv2/IPsec route-based + NAT-T       & Baseline; DMVPN a evaluar. \\
Starlink IPv4   & CGNAT \texttt{100.64.0.0/10} (pública opcional) & Registrar el observado (G9). \\
SFP             & Solo \textbf{1\,Gb/s SFP}             & \textbf{Prohibido SFP+}; P/N validado en G6. \\
\bottomrule
\end{tabularx}

\vspace{1em}
\begin{nota}
\small
Todo lo anterior se materializa siguiendo \file{04-pasos-montaje.md}.
Los supuestos marcados (G*) se cierran según \file{07-riesgos-y-verificaciones.md}
\textbf{antes} de depender de ellos.
\end{nota}

\end{document}
